\Chapter{Organizational assumptions for EMR system implementation}
\label{chapter:implementation}

\section{Introduction}

The implementation of the EMR system at Szpital im. Bajgla is not only a technological project, but also a major organizational change covering the way medical and administrative staff work every day. In addition to preparing the system itself, it is also necessary to plan the implementation method, documentation migration and the organization of users' work during the transition period.

The purpose of the presented assumptions is to reduce organizational and technical problems related to the facility's digitalization while maintaining continuity of the hospital's operation. These assumptions are also intended to reduce the risk of errors during implementation and to facilitate staff adaptation to the new working environment.

\section{Phased system implementation}

Due to the scale of Szpital im. Bajgla's activity, the system implementation should be carried out in stages. This solution allows reducing organizational risk and avoiding a situation in which a simultaneous change in the working method of the entire facility could lead to problems with the functioning of departments.

In the first stage, the system should be implemented in selected departments with a lower degree of organizational complexity. This will allow testing the operation of basic functionalities and introducing possible corrections before expanding the system to the remaining hospital units. Only after completing the pilot phase could the solution be implemented in departments requiring rapid information flow, such as the Emergency Department or the intensive care unit.

The migration of paper documentation should also be carried out gradually. In the first place, current patient data and documentation used in the ongoing work of departments would be transferred to the system. Archival documentation could be digitized in stages at a later period, which would allow reducing costs and the organizational burden associated with simultaneously processing a very large amount of data.

During the transition period, the simultaneous operation of some paper and digital documentation is provided for. Such a solution enables gradual familiarization of staff with the new system and reduces the risk of losing access to data during the initial implementation stage.

\section{Staff training and work organization}

Implementing the new system requires appropriate preparation of medical and administrative staff to work in a digital environment. For this reason, before starting to use the system, it is necessary to conduct training covering basic functionalities and procedures for working with electronic documentation.

Training should be adapted to the scope of duties of individual user groups. Doctors, nurses, laboratory technicians and administrative employees use different elements of the system, which is why the scope of information provided should correspond to the actual needs of a given group.

In order to reduce organizational problems, the preparation of basic procedures for using the system is also provided for. This concerns, among other things, the method of adding documentation, updating patient data, transferring information between departments and proceeding in emergency situations.

An important element of work organization is also determining responsibility for entering and updating data. Each user should have an individual account enabling identification of actions performed in the system. This solution allows reducing the risk of errors and increases the security of medical documentation.

\section{Technical assumptions and system security}

The designed system assumes the use of a central database enabling the storage of patient documentation in a single coherent environment. Thanks to this, it will be possible to synchronize data between departments and access current information in real time.

An important project assumption is to ensure an appropriate level of medical data security. The system should use individual user accounts, a division of access levels and encryption of data and information transmission. Thanks to this, it will be possible to limit access to documentation only to persons with appropriate permissions.

The project also provides for the creation of regular backups and procedures for recovering data after a system failure. This solution is intended to reduce the risk of losing medical documentation and to enable rapid restoration of system operation in the event of technical problems.

In addition, the system should enable monitoring of user activity and recording of changes made to medical documentation. This functionality is important both from the perspective of data security and compliance with the requirements concerning the protection of medical information in force in the European Union.

\section{Procedure in emergency situations}

Due to the round-the-clock nature of the hospital's work, it is necessary to prepare procedures enabling the facility to function also in the event of temporary unavailability of the system. In the event of a failure, staff should have the possibility of switching to a simplified working mode and temporarily using substitute documentation until full system functionality is restored.

The project also provides for permanent technical support after system implementation. In the initial period of the solution's operation, rapid response to problems reported by users and the possibility of ongoing organizational and technical corrections will be particularly important.
